%%
%% NeurIPS 2025 LaTeX Template
%% Template for Conference on Neural Information Processing Systems submissions
%%

\documentclass{article}

% NeurIPS Packages
\usepackage{neurips_2025}  % NeurIPS style file (needs to be downloaded separately)
\usepackage[T1]{fontenc}
\usepackage{hyperref}
\usepackage{url}
\usepackage{booktabs}
\usepackage{amsfonts}
\usepackage{amsmath}
\usepackage{amssymb}
\usepackage{graphicx}
\usepackage{algorithm}
\usepackage{algorithmic}
\usepackage{xcolor}
\usepackage{subcaption}

% For hyperlinks
\hypersetup{
    colorlinks=true,
    linkcolor=blue,
    citecolor=blue,
    urlcolor=blue
}

% Math operators
\DeclareMathOperator*{\argmin}{arg\,min}
\DeclareMathOperator*{\argmax}{arg\,max}

% Title
\title{Your Paper Title Here}

% Authors
\author{%
  First Author\thanks{Equal contribution.} \And
  Second Author \And
  Third Author \\
  \texttt{first.author@university.edu} \\
  \texttt{second.author@university.edu} \\
  \texttt{third.author@university.edu} \\
  \\
  Department of Computer Science\\
  University Name\\
}

\begin{document}

\maketitle

% Abstract
\begin{abstract}
Your abstract should briefly summarize the main contributions of your paper.
Keep it concise (typically 150-250 words). Cover: problem, approach, results, implications.
Avoid citations and complex mathematical notation here.
\end{abstract}

% Main content
\section{Introduction}
\label{sec:introduction}

Introduce the problem you address and motivate your work:
\begin{itemize}
    \item What is the challenge?
    \item Why is it important?
    \item What gap does your work fill?
\end{itemize}

State your contributions clearly:

\begin{enumerate}
    \item We propose a novel method for [problem/task]
    \item We provide theoretical analysis showing [key insight]
    \item We demonstrate empirically that our method achieves [results]
\end{enumerate}

\section{Related Work}
\label{sec:related}

Survey relevant prior work and position your contribution:
\begin{itemize}
    \item Category 1 of related work \citep{author2024paper}
    \item Category 2 of related work
    \item How your work differs from existing approaches
\end{itemize}

\section{Background}
\label{sec:background}

\subsection{Problem Setup}

Define the problem formally. Let $\mathcal{X} \subset \mathbb{R}^d$ be the input space
and $\mathcal{Y}$ be the label space. Given training data
$\mathcal{D} = \{(x_i, y_i)\}_{i=1}^n$, we aim to learn a mapping
$f_\theta: \mathcal{X} \rightarrow \mathcal{Y}$.

\subsection{Notation}

Use consistent notation throughout:
\begin{itemize}
    \item $x \in \mathcal{X}$: input
    \item $y \in \mathcal{Y}$: output/label
    \item $\theta$: model parameters
    \item $\mathcal{L}$: loss function
\end{itemize}

\section{Method}
\label{sec:method}

\subsection{Overview}

Describe your approach at a high level.

\subsection{Proposed Approach}

Detail the technical contributions:

\textbf{Key Idea 1.} Describe the first key insight.

\textbf{Key Idea 2.} Describe the second key insight.

\subsection{Objective Function}

Define your objective:

\begin{equation}
    \min_{\theta} \mathcal{L}(\theta) = \mathbb{E}_{(x,y) \sim \mathcal{D}}
    \left[ \ell(f_\theta(x), y) \right] + \lambda R(\theta)
    \label{eq:objective}
\end{equation}

where $\ell$ is the task loss, $R(\theta)$ is a regularization term,
and $\lambda > 0$ is a hyperparameter.

\subsection{Optimization}

\begin{algorithm}[tb]
   \caption{Proposed Algorithm}
   \label{alg:algorithm}
\begin{algorithmic}
   \STATE \textbf{Input:} Data $\mathcal{D}$, learning rate $\eta$, iterations $T$
   \STATE \textbf{Initialize:} $\theta_0$
   \FOR{$t = 1$ to $T$}
       \STATE Sample batch $\mathcal{B} \sim \mathcal{D}$
       \STATE Compute loss: $L_t = \frac{1}{|\mathcal{B}|}\sum_{(x,y)\in\mathcal{B}} \ell(f_\theta(x), y)$
       \STATE Update: $\theta_{t+1} = \theta_t - \eta \nabla_\theta L_t$
   \ENDFOR
   \STATE \textbf{Return:} $\theta_T$
\end{algorithmic}
\end{algorithm}

\subsection{Theoretical Analysis}

If applicable, include theoretical results:

\begin{theorem}
\label{thm:main}
Under conditions [A1-A3], Algorithm~\ref{alg:algorithm} converges to a
stationary point at rate $O(1/\sqrt{T})$.
\end{theorem}

\begin{proof}
(Proof sketch or reference to appendix)
\end{proof}

\section{Experiments}
\label{sec:experiments}

\subsection{Experimental Setup}

\textbf{Datasets.} Describe datasets used:
\begin{itemize}
    \item Dataset 1: Description and statistics
    \item Dataset 2: Description and statistics
\end{itemize}

\textbf{Baselines.} Compare against:
\begin{itemize}
    \item Baseline method 1 \citep{ref1}
    \item Baseline method 2 \citep{ref2}
\end{itemize}

\textbf{Implementation.} Hyperparameters and training details.

\subsection{Main Results}

Present quantitative results in tables:

\begin{table}[t]
    \centering
    \caption{Main results on benchmark datasets.}
    \label{tab:results}
    \begin{tabular}{lccc}
        \toprule
        Method & Metric 1 & Metric 2 & Metric 3 \\
        \midrule
        Baseline 1 & 85.2 & 0.84 & 92.1 \\
        Baseline 2 & 87.3 & 0.86 & 93.5 \\
        \textbf{Ours} & \textbf{91.5} & \textbf{0.90} & \textbf{96.2} \\
        \bottomrule
    \end{tabular}
\end{table}

\subsection{Ablation Study}

Analyze the contribution of each component:

\begin{table}[t]
    \centering
    \caption{Ablation study results.}
    \label{tab:ablation}
    \begin{tabular}{lc}
        \toprule
        Variant & Performance \\
        \midrule
        Full model & 91.5 \\
        w/o component 1 & 88.2 \\
        w/o component 2 & 89.7 \\
        w/o component 3 & 90.1 \\
        \bottomrule
    \end{tabular}
\end{table}

\subsection{Analysis}

Discuss results, failure cases, and insights.

\section{Conclusion}
\label{sec:conclusion}

Summarize key contributions and future directions.

\paragraph{Limitations.} Acknowledge limitations of your approach.

\paragraph{Future Work.} Suggest extensions and improvements.

\paragraph{Broader Impact.} Discuss societal implications if relevant.

\section*{Acknowledgments}

Thank funding agencies, collaborators, and reviewers.

\paragraph{Funding Disclosure} This work was supported by [grants].

% Bibliography
\bibliographystyle{plain}
\bibliography{references}

% Appendix
\appendix

\section{Proof of Theorem~\ref{thm:main}}
\label{app:proof}

Provide detailed proofs here.

\section{Additional Experimental Details}
\label{app:experiments}

Include:
\begin{itemize}
    \item Full hyperparameter tables
    \item Additional results and visualizations
    \item Implementation details
\end{itemize}

\end{document}
